% Unit 086 terjemahan Bahasa Indonesia dari chapter6.tex baris 2604--2927.
% Sumber: Methods of Algebra, Volume 2, commit
% 9a5803ff77dd3257484cb177f851a73770a59dd3 (CC BY 4.0).
% stable-unit-id: o014.aljabr2.chapter6.nonabelian-cohomology
% Cakupan: Bagian 6.13 lengkap, ``Kohomologi Nonabelian'';
% berhenti tepat sebelum lingkungan Exercises pada baris 2928.

% segment-id: o014.li2.chapter6.nonabelian-cohomology.h001
\section{Kohomologi Nonabelian}\label{sec:nonabelian-coh}
\hypertarget{o014.aljabr2.chapter6.nonabelian-cohomology}{ }
% segment-id: o014.li2.chapter6.nonabelian-cohomology.p001
Misalkan $G$ grup. Kohomologi $\Hm^n(G,M)$ yang dibahas sebelumnya selalu
mensyaratkan ``koefisien'' $M$ memiliki struktur modul, atau setidaknya
merupakan modul-$\Z$, yakni grup abelian. Bagian ini bertujuan membahas kasus
ketika koefisiennya grup nonabelian. Mulai sekarang, kita memakai simbol $A$
sebagai pengganti $M$.

% segment-id: o014.li2.chapter6.nonabelian-cohomology.d001
\begin{definisi}\label{def:G-group}
	\index[sym1]{G-Grp@$G\dcate{Grp}$}
	\index{grup-G@grup-$G$}
	Misalkan $A$ grup, dengan operasi biner ditulis secara multiplikatif dan
	unsur identitas ditulis $1$. Andaikan $A$ dilengkapi aksi kiri grup $G$,
	yang ditulis
% segment-id: o014.li2.chapter6.nonabelian-cohomology.q001
	\[ G \times A \to A, \quad (g, a) \mapsto {}^g a . \]
	Jika aksi grup tersebut kompatibel dengan struktur perkalian $A$,
% segment-id: o014.li2.chapter6.nonabelian-cohomology.q002
	\[ {}^g (a_1 a_2) = {}^g a_1 {}^g a_2, \quad {}^g 1 = 1, \]
	maka $A$ disebut grup dengan aksi $G$, atau singkatnya \emph{grup-$G$}.
	Homomorfisme antargrup-$G$ didefinisikan sebagai homomorfisme grup
	yang mempertahankan aksi $G$. Kategori yang dibentuknya dituliskan
	$G\dcate{Grp}$.
\end{definisi}

% segment-id: o014.li2.chapter6.nonabelian-cohomology.p002
Sebentar lagi akan terlihat bahwa, untuk grup-$G$ $A$, tanpa struktur tambahan
atau komutativitas, hanya $\Hm^0$ dan $\Hm^1$ yang dapat didefinisikan secara
wajar dan berguna. Dalam keadaan ini, $\Hm^1$ belum tentu grup, melainkan
himpunan bertitik.

% segment-id: o014.li2.chapter6.nonabelian-cohomology.d002
\begin{definisi}\label{def:pointed-set}
	\index{himpunan bertitik (pointed set)}
	\index[sym1]{Set-pt@$\cate{Set}_\bullet$, $G\dcate{Set}_\bullet$}
	Suatu \emph{himpunan bertitik} adalah data $(X,x_0)$, dengan $X$ himpunan
	dan $x_0\in X$ (disebut titik pangkal). Morfisme dari $(X,x_0)$ ke
	$(Y,y_0)$ adalah pemetaan $f:X\to Y$ yang memenuhi $f(x_0)=y_0$. Data ini
	membentuk kategori $\cate{Set}_\bullet$.
	
	Misalkan $(X,x_0)$ himpunan bertitik. Jika grup $G$ beraksi dari kiri pada
	$X$, ditulis $(g,x)\mapsto{}^gx$, dan ${}^gx_0=x_0$ selalu berlaku, kita
	katakan bahwa $G$ beraksi pada $(X,x_0)$. Dalam hal ini, dapat didefinisikan
	subhimpunan invarian bertitik
% segment-id: o014.li2.chapter6.nonabelian-cohomology.q003
	\[ (X, x_0)^G := \left\{ x \in X: \forall g \in G, \; gx = x \right\}, \quad \text{titik pangkal} = x_0. \]
	Semua himpunan bertitik dengan aksi $G$, beserta morfisme yang kompatibel
	dengan aksi $G$, membentuk kategori yang dituliskan
	$G\dcate{Set}_\bullet$.
\end{definisi}

% segment-id: o014.li2.chapter6.nonabelian-cohomology.p003
Jika tidak menimbulkan kerancuan, objek $(X,x_0)$ dari
$\cate{Set}_\bullet$ sering cukup dituliskan sebagai $X$, dan dalam
$G\dcate{Set}_\bullet$, $(X,x_0)^G$ dituliskan sebagai $X^G$.

% segment-id: o014.li2.chapter6.nonabelian-cohomology.p004
Setiap grup tentu menjadi himpunan bertitik dengan memilih unsur identitas
sebagai titik pangkal. Ini memberikan funktor yang jelas
$\cate{Grp}\to\cate{Set}_\bullet$; demikian pula terdapat funktor
$G\dcate{Grp}\to G\dcate{Set}_\bullet$. Perhatikan bahwa subhimpunan
invarian $A^G$ dari grup-$G$ $A$ merupakan subgrupnya.

% segment-id: o014.li2.chapter6.nonabelian-cohomology.d003
\begin{definisi}\label{def:pointed-exact}
	\index{barisan eksak}
	Misalkan
	$(X,x_0)\xrightarrow{f}(Y,y_0)\xrightarrow{g}(Z,z_0)$ morfisme-morfisme
	dalam $\cate{Set}_\bullet$ (atau $G\dcate{Set}_\bullet$). Barisan morfisme
	ini disebut eksak jika $g^{-1}(z_0)=\Image(f)$. Dengan cara ini,
	barisan eksak dalam $\cate{Set}_\bullet$ (atau
	$G\dcate{Set}_\bullet$) dapat didefinisikan.
\end{definisi}

% segment-id: o014.li2.chapter6.nonabelian-cohomology.p005
Untuk grup abelian, definisi ini kembali kepada keeksakan yang lazim.
Sekarang $\Hm^0$ dan $\Hm^1$ dapat didefinisikan bagi setiap grup-$G$ $A$.

% segment-id: o014.li2.chapter6.nonabelian-cohomology.d004
\begin{definisi}\label{def:nonabelian-H}
	\index{kohomologi grup!kasus nonabelian}
	Misalkan $A$ objek $G\dcate{Set}_\bullet$. Definisikan himpunan bertitik
% segment-id: o014.li2.chapter6.nonabelian-cohomology.q004
	\[ \Hm^0(G, A) = Z^0(G, A) := A^G. \]
	
	Jika $A$ grup-$G$, maka $\Hm^0(G,A)$ merupakan grup. Dalam hal ini,
	definisikan lebih lanjut
% segment-id: o014.li2.chapter6.nonabelian-cohomology.q005
	\begin{align*}
		Z^1(G, A) & := \left\{\begin{array}{r|l}
			c: G \to A & \forall g_1, g_2 \in G, \\
			& c(g_1 g_2) = c(g_1) \cdot {}^{g_1} c(g_2)
		\end{array}\right\}, \\
		\Hm^1(G, A) & := Z^1(G, A) / \sim ,
	\end{align*}
	dengan relasi ekuivalensi $c\sim c'$ berarti bahwa terdapat $a\in A$
	sedemikian sehingga, untuk setiap $g\in G$,
% segment-id: o014.li2.chapter6.nonabelian-cohomology.q006
	\[ c'(g) = a^{-1} \cdot c(g) \cdot {}^g a . \]
	Mudah dilihat bahwa ini memang relasi ekuivalensi—relasi tersebut berasal
	dari suatu aksi kanan $A$—dan $\Hm^1(G,A)$ secara alami menjadi himpunan
	bertitik. Titik pangkalnya ialah kelas ekuivalensi yang ditentukan oleh
	pemetaan konstan $c(g)=1$.
\end{definisi}

% segment-id: o014.li2.chapter6.nonabelian-cohomology.p006
Perhatikan bahwa memasukkan $g_1=g_2=1_G$ ke syarat dalam $Z^1(G,A)$
memberikan $c(1_G)=1$. Selanjutnya, $[c]$ menyatakan kelas ekuivalensi dari
$c\in Z^1(G,A)$.

% segment-id: o014.li2.chapter6.nonabelian-cohomology.r001
\begin{catatan}[$\Hm^1$ sebagai kelas konjugasi penampang]\label{rem:Z1-vs-section}
	Dari grup-$G$ $A$, bentuk produk semilangsung $A\rtimes G$, dan tuliskan
	$\pi$ untuk homomorfisme proyeksi $A\rtimes G\to G$. Unsur-unsur
	$Z^1(G,A)$ berkorespondensi secara bijektif dengan homomorfisme grup
	$\sigma:G\to A\rtimes G$ yang memenuhi $\pi\sigma=\identity_G$ (yakni
	``penampang'' dari $\pi$), melalui pemetaan
	$c\in Z^1(G,A)$ ke $\sigma(g)=(c(g),g)$. Lebih lanjut, $[c']=[c]$ jika
	dan hanya jika terdapat $a\in A$ sedemikian sehingga
	$\sigma'=a^{-1}\sigma a$, dengan $A$ disertakan sebagai subgrup
	$A\rtimes G$ dan konjugasi dilakukan titik demi titik.
\end{catatan}

% segment-id: o014.li2.chapter6.nonabelian-cohomology.p007
Jika $A$ grup abelian, konstruksi ini kembali ke deskripsi $\Hm^0$ dan
$\Hm^1$ melalui kompleks kokrantai standar
(Proposisi~\ref{prop:groupcoh-cochain}). Dalam kasus nonabelian,
konstruksi tersebut tetap funktorial, sebagai berikut.
% segment-id: o014.li2.chapter6.nonabelian-cohomology.l001
\begin{itemize}
% segment-id: o014.li2.chapter6.nonabelian-cohomology.i001
	\item Setiap morfisme $f:A_1\to A_2$ menginduksi
% segment-id: o014.li2.chapter6.nonabelian-cohomology.q007
	\begin{gather*}
		\Hm^0(f): \Hm^0(G, A_1) \to \Hm^0(G, A_2), \quad
		\Hm^1(f): \Hm^1(G, A_1) \to \Hm^1(G, A_2);
	\end{gather*}
	Pemetaan pertama jelas, sedangkan pemetaan kedua diberikan oleh
	$\Hm^1(f)([c])=[f\circ c]$.
	
% segment-id: o014.li2.chapter6.nonabelian-cohomology.i002
	\item Misalkan $\varphi:H\to G$ homomorfisme grup. Tarik kembali aksi $G$
	pada grup-$G$ $A$ menjadi aksi $H$, sehingga diperoleh grup-$H$
	$\varphi^*A$.
% segment-id: o014.li2.chapter6.nonabelian-cohomology.l002
	\begin{compactitem}
% segment-id: o014.li2.chapter6.nonabelian-cohomology.i003
		\item Definisikan morfisme kanonik
		$\Hm^0(G,A)\to\Hm^0(H,\varphi^*A)$ sebagai inklusi yang jelas
		$A^G\hookrightarrow A^{\varphi(H)}=(\varphi^*A)^H$.
% segment-id: o014.li2.chapter6.nonabelian-cohomology.i004
		\item Definisikan morfisme kanonik himpunan bertitik
% segment-id: o014.li2.chapter6.nonabelian-cohomology.q008
		\[ \Hm^1(G, A) \to \Hm^1(H, \varphi^* A). \]
		melalui $[c]\mapsto[c\circ\varphi]$.
	\end{compactitem}
\end{itemize}

% segment-id: o014.li2.chapter6.nonabelian-cohomology.p008
Pemetaan-pemetaan kanonik ini memiliki berbagai sifat seperti yang terlihat
dalam \S\ref{sec:change-of-group}. Semuanya juga dapat dirangkum sebagai
berikut. Untuk $\varphi:H\to G$, misalkan $A_1$ objek
$G\dcate{Set}_\bullet$ (atau $G\dcate{Grp}$), $A_2$ objek
$H\dcate{Set}_\bullet$ (atau $H\dcate{Grp}$), dan morfisme $f:A_1\to A_2$
ekuivarian terhadap $\varphi$ seperti dalam
Definisi~\ref{def:group-equivariance}. Terdapat pemetaan terkait
$\Hm^i(f):\Hm^i(G,A_1)\to\Hm^i(H,A_2)$, sehingga analog
Proposisi~\ref{prop:grp-equiv} berbentuk:
% segment-id: o014.li2.chapter6.nonabelian-cohomology.l003
\begin{itemize}
% segment-id: o014.li2.chapter6.nonabelian-cohomology.i005
	\item $\Hm^0(f)$ ialah restriksi $f$ pada $A_1^G\to A_2^H$;
% segment-id: o014.li2.chapter6.nonabelian-cohomology.i006
	\item $\Hm^1(f)$ memetakan $[c]$ ke $[f\circ c\circ\varphi]$.
\end{itemize}

% segment-id: o014.li2.chapter6.nonabelian-cohomology.c001
\begin{konvensi}
	Tuliskan $1$ untuk objek nol dalam $\cate{Set}_\bullet$ atau
	$G\dcate{Set}_\bullet$ yang ditentukan oleh himpunan satu unsur.
\end{konvensi}

% segment-id: o014.li2.chapter6.nonabelian-cohomology.p009
Sekarang kita masuk ke pokok bahasan. Selanjutnya, tinjau barisan eksak pendek
dalam $G\dcate{Set}_\bullet$ (Definisi~\ref{def:pointed-exact})
% segment-id: o014.li2.chapter6.nonabelian-cohomology.q009
\[ 1 \to A \xrightarrow{u} B \xrightarrow{v} C \to 1, \]
dengan $A$ dan $B$ diandaikan grup-$G$, serta $u$ homomorfisme grup.
Keeksakan dengan demikian menyiratkan bahwa $u$ injektif dan $v$ surjektif.
Kita mensyaratkan lebih lanjut:
% segment-id: o014.li2.chapter6.nonabelian-cohomology.q010
\begin{equation}\label{eqn:pointed-ses-condition}
	\begin{gathered}
		\forall b, b' \in B, \; \left[ v(b) = v(b') \iff \exists a \in A, \; b' = b u(a) \right], \\
		\text{dengan kata lain: $C$ diidentifikasi dengan $B/u(A)$, titik pangkal $= 1 \cdot u(A)$.}
	\end{gathered}
\end{equation}

% segment-id: o014.li2.chapter6.nonabelian-cohomology.p010
Jika $C$ grup dan $v$ homomorfisme grup, dengan meninjau
$b^{-1}b'\in v^{-1}(1)=\Image(u)$ terlihat bahwa
\eqref{eqn:pointed-ses-condition} berlaku otomatis, dan dalam keadaan ini
$u(A)\lhd B$. Namun, banyak penerapan tidak berada dalam keadaan ini.

% segment-id: o014.li2.chapter6.nonabelian-cohomology.p011
Tujuan kita sekarang ialah membangun barisan eksak yang ``sepanjang mungkin''
beserta morfisme penghubung terkait dari data ini. Tanpa penjelasan lebih
lanjut, kita identifikasikan $A$ dengan subgrup $B$ dan menghilangkan simbol
$u$.

% segment-id: o014.li2.chapter6.nonabelian-cohomology.p012
\textbf{Kasus derajat nol.}\;
Pertama, definisikan morfisme penghubung
$\delta^0:\Hm^0(G,C)\to\Hm^1(G,A)$. Misalkan $c\in C^G$. Pilih sembarang
$b\in v^{-1}(c)$. Karena $v$ mempertahankan aksi $G$, dari
\eqref{eqn:pointed-ses-condition} diketahui bahwa, untuk setiap $g\in G$,
terdapat tepat satu $a(g)\in A$ sedemikian sehingga ${}^gb=ba(g)$. Maka
$g\mapsto a(g)$ memberikan unsur $a$ dari $Z^1(G,A)$, sebab
% segment-id: o014.li2.chapter6.nonabelian-cohomology.q011
\[ a(g_1 g_2) = b^{-1} \cdot {}^{g_1 g_2} b = b^{-1} \cdot {}^{g_1} b \cdot {}^{g_1} (b^{-1} \cdot {}^{g_2} b) = a(g) \cdot {}^{g_1} a(g_2). \]
Jika $b',b\in v^{-1}(c)$, terdapat $\alpha\in A$ sedemikian sehingga
$b'=b\alpha$, dan karena itu
% segment-id: o014.li2.chapter6.nonabelian-cohomology.q012
\[ (b')^{-1} \cdot {}^g b' = \alpha^{-1} \cdot b^{-1} \cdot {}^g b \cdot {}^g \alpha. \]
Ini menunjukkan bahwa kelas $[a]\in\Hm^1(G,A)$ hanya ditentukan oleh $c$.
Jika $c$ merupakan titik pangkal $C$, kita dapat mengambil $b=1\in B$.
Dengan demikian diperoleh morfisme dalam $\cate{Set}_\bullet$
$\delta^0:\Hm^0(G,C)\to\Hm^1(G,A)$.

% segment-id: o014.li2.chapter6.nonabelian-cohomology.p013
\textbf{Kasus derajat nol: subgrup pusat.}\;
Andaikan $C$ grup, $v$ homomorfisme grup, dan $A$ termuat dalam pusat $Z_B$
dari grup $B$; khususnya, $A$ grup abelian. Dalam keadaan ini,
$\Hm^0(G,C)=C^G$ dan $\Hm^1(G,A)$ keduanya grup. Dari definisi $\delta^0$
dan syarat $A\subset Z_B$, mudah diperiksa bahwa $\delta^0$ merupakan
homomorfisme grup.

% segment-id: o014.li2.chapter6.nonabelian-cohomology.p014
\textbf{Kasus derajat satu.}\;
Tetap andaikan $C$ grup, $v$ homomorfisme grup, dan $A\subset Z_B$; maka
$\Hm^2(G,A)$ dapat dideskripsikan melalui kompleks kokrantai standar. Kita
akan mendefinisikan $\delta^1:\Hm^1(G,C)\to\Hm^2(G,A)$. Misalkan
$c\in Z^1(G,C)$. Untuk setiap $g\in G$, pilih $b(g)\in v^{-1}(c(g))$.
Karena \eqref{eqn:pointed-ses-condition} selalu berlaku dan $v$
mempertahankan aksi $G$, terdapat tepat satu pemetaan $a:G^2\to A$ yang
memenuhi
% segment-id: o014.li2.chapter6.nonabelian-cohomology.q013
\[ b(g_1) \cdot {}^{g_1} b(g_2) = a(g_1, g_2) b(g_1 g_2). \]

% segment-id: o014.li2.chapter6.nonabelian-cohomology.p015
Sekarang periksa bahwa $a\in Z^2(G,A)$. Ini ekuivalen dengan
% segment-id: o014.li2.chapter6.nonabelian-cohomology.q014
\[ a(g_1, g_2)^{-1} \cdot a(g_1, g_2 g_3) \cdot a(g_1 g_2, g_3)^{-1} \cdot {}^{g_1} a(g_2, g_3) = 1. \]
Uraikan ruas kiri sebagai ekspresi dalam $B$:
% segment-id: o014.li2.chapter6.nonabelian-cohomology.q015
\begin{multline*}
	\underbracket{b(g_1 g_2) \cdot {}^{g_1} b(g_2)^{-1} \cdot b(g_1)^{-1}}_{a(g_1, g_2)^{-1}} \cdot \underbracket{b(g_1) \cdot {}^{g_1} b(g_2 g_3) b(g_1 g_2 g_3)^{-1}}_{a(g_1, g_2 g_3)} \\
	\underbracket{b(g_1 g_2 g_3) \cdot {}^{g_1 g_2} b(g_3)^{-1} \cdot b(g_1 g_2)^{-1}}_{a(g_1 g_2, g_3)^{-1}} \; \cdot \; {}^{g_1} a(g_2, g_3),
\end{multline*}
lalu gunakan $A\subset Z_B$ untuk menyederhanakannya menjadi
% segment-id: o014.li2.chapter6.nonabelian-cohomology.q016
\begin{multline*}
	b(g_1 g_2) \cdot {}^{g_1} b(g_2)^{-1} \cdot {}^{g_1} b(g_2 g_3) \cdot {}^{g_1 g_2} b(g_3)^{-1}  b(g_1 g_2)^{-1} \cdot {}^{g_1} a(g_2, g_3) \\
	= b(g_1 g_2) \cdot {}^{g_1} b(g_2)^{-1} \cdot \underbracket{{}^{g_1} b(g_2) \cdot {}^{g_1 g_2} b(g_3) \cdot {}^{g_1} b(g_2 g_3)^{-1}}_{{}^{g_1} a(g_2, g_3)} \cdot {}^{g_1} b(g_2 g_3) \cdot {}^{g_1 g_2} b(g_3)^{-1}  b(g_1 g_2)^{-1} .
\end{multline*}
Ekspresi terakhir jelas sama dengan $1$, sehingga $a\in Z^2(G,A)$.

% segment-id: o014.li2.chapter6.nonabelian-cohomology.p016
Jika pilihan setiap $b(g)$ diubah dan dituliskan
$b'(g)=\alpha(g)b(g)$, dengan $\alpha:G\to A$ sembarang pemetaan, maka
$a'(g_1,g_2)$ yang bersesuaian adalah
% segment-id: o014.li2.chapter6.nonabelian-cohomology.q017
\begin{align*}
	a'(g_1, g_2) & = \alpha(g_1) \cdot b(g_1) \cdot {}^{g_1} \alpha(g_2) \cdot {}^{g_1} b(g_2) b(g_1 g_2)^{-1} \alpha(g_1 g_2)^{-1} \\
	& = \alpha(g_1) \cdot {}^{g_1} \alpha(g_2) \cdot \alpha(g_1 g_2)^{-1} \cdot a(g_1, g_2) \quad (\because\; A \subset Z_B).
\end{align*}
Dengan demikian, $a,a'\in Z^2(G,A)$ hanya berbeda dengan faktor berupa suatu
unsur $B^2(G,A)$.

% segment-id: o014.li2.chapter6.nonabelian-cohomology.p017
Terakhir, tinjau kasus ketika $c$ diganti dengan
$c'(g)=\gamma^{-1}\cdot c(g)\cdot{}^g\gamma$, dengan $\gamma\in C$. Pilih
sembarang $\beta\in v^{-1}(\gamma)$. Kita dapat mengambil
$b'(g)=\beta^{-1}\cdot b(g)\cdot{}^g\beta$ untuk $b'$ yang bersesuaian
dengan $c'$. Pemetaan $a':G^2\to A$ yang dibangun dari $b'$ berbentuk
% segment-id: o014.li2.chapter6.nonabelian-cohomology.q018
\begin{align*}
	a'(g_1, g_2) & = \left( \beta^{-1} \cdot b(g_1) \cdot {}^{g_1} \beta \right) \cdot {}^{g_1} \left( \beta^{-1} \cdot b(g_2) \cdot {}^{g_2} \beta \right) \cdot \left( {}^{g_1 g_2} \beta^{-1} \cdot b(g_1 g_2)^{-1} \cdot \beta \right) \\
	& = \beta^{-1} \cdot \underbracket{b(g_1) \cdot {}^{g_1} b(g_2) \cdot b(g_1 g_2)^{-1}}_{= a(g_1, g_2)} \cdot \beta \\
	& = a(g_1, g_2).
\end{align*}

% segment-id: o014.li2.chapter6.nonabelian-cohomology.p018
Jadi kelas $a$ dalam $\Hm^2(G,A)$ ditentukan oleh
$[c]\in\Hm^1(G,C)$. Jika $c$ adalah pemetaan konstan bernilai $1$, $b$ juga
dapat dipilih konstan bernilai $1$, dan demikian pula $a$. Dengan demikian,
diperoleh morfisme $\cate{Set}_\bullet$
$\delta^1:\Hm^1(G,C)\to\Hm^2(G,A)$.

% segment-id: o014.li2.chapter6.nonabelian-cohomology.th001
\begin{teorema}\label{prop:nonabelian-long}
	Misalkan $1\to A\xrightarrow{u}B\xrightarrow{v}C\to1$ barisan eksak
	pendek dalam $G\dcate{Set}_\bullet$, dengan $A$ dan $B$ grup-$G$ serta
	$u$ homomorfisme grup.
% segment-id: o014.li2.chapter6.nonabelian-cohomology.l004
	\begin{enumerate}[(i)]
% segment-id: o014.li2.chapter6.nonabelian-cohomology.i007
		\item Jika syarat \eqref{eqn:pointed-ses-condition} berlaku, terdapat
		barisan eksak dalam $\cate{Set}_\bullet$
% segment-id: o014.li2.chapter6.nonabelian-cohomology.g001
		\begin{equation*}\begin{tikzcd}[row sep=large]
			1 \arrow[r] & \Hm^0(G, A) \arrow[r, "{\Hm^0(u)}"] & \Hm^0(G, B) \arrow[r, "{\Hm^0(v)}"] \arrow[phantom, d, ""{coordinate, name=A}] & \Hm^0(G, C) \arrow[dll, rounded corners, "{\delta^0}" description, to path={
				-- ([xshift=2ex]\tikztostart.east)
				|- (A) [near end]\tikztonodes
				-| ([xshift=-2ex]\tikztotarget.west)
				-- (\tikztotarget)}] \\
			& \Hm^1(G, A) \arrow[r, "{\Hm^1(u)}"] & \Hm^1(G, B). &
		\end{tikzcd}\end{equation*}
% segment-id: o014.li2.chapter6.nonabelian-cohomology.i008
		\item Selain asumsi di atas, andaikan $C$ juga grup dan $v$
		homomorfisme grup; hal ini otomatis menyiratkan
		\eqref{eqn:pointed-ses-condition}. Barisan eksak tersebut memanjang
		menjadi
% segment-id: o014.li2.chapter6.nonabelian-cohomology.g002
		\begin{equation*}\begin{tikzcd}[row sep=large]
			1 \arrow[r] & \Hm^0(G, A) \arrow[r, "{\Hm^0(u)}"] & \Hm^0(G, B) \arrow[r, "{\Hm^0(v)}"] \arrow[phantom, d, ""{coordinate, name=A}] & \Hm^0(G, C) \arrow[dll, rounded corners, "{\delta^0}" description, to path={
				-- ([xshift=2ex]\tikztostart.east)
				|- (A) [near end]\tikztonodes
				-| ([xshift=-2ex]\tikztotarget.west)
				-- (\tikztotarget)}] \\
			& \Hm^1(G, A) \arrow[r, "{\Hm^1(u)}"] & \Hm^1(G, B) \arrow[r, "{\Hm^1(v)}"] & \Hm^1(G, C).
		\end{tikzcd}\end{equation*}
% segment-id: o014.li2.chapter6.nonabelian-cohomology.i009
		\item Jika disyaratkan lebih lanjut bahwa $u(A)$ termuat dalam pusat
		$Z_B$ dari $B$, maka $\delta^0$ merupakan homomorfisme grup dan barisan
		eksak tersebut memanjang menjadi
% segment-id: o014.li2.chapter6.nonabelian-cohomology.g003
		\begin{equation*}\begin{tikzcd}[row sep=large]
			1 \arrow[r] & \Hm^0(G, A) \arrow[r, "{\Hm^0(u)}"] & \Hm^0(G, B) \arrow[r, "{\Hm^0(v)}"] \arrow[phantom, d, ""{coordinate, name=A}] & \Hm^0(G, C) \arrow[dll, rounded corners, "{\delta^0}" description, to path={
				-- ([xshift=2ex]\tikztostart.east)
				|- (A) [near end]\tikztonodes
				-| ([xshift=-2ex]\tikztotarget.west)
				-- (\tikztotarget)}] \\
			& \Hm^1(G, A) \arrow[r, "{\Hm^1(u)}"] & \Hm^1(G, B) \arrow[r, "{\Hm^1(v)}"] \arrow[phantom, d, ""{coordinate, name=B}] & \Hm^1(G, C) \arrow[dll, rounded corners, "{\delta^1}" description, to path={
				-- ([xshift=2ex]\tikztostart.east)
				|- (B) [near end]\tikztonodes
				-| ([xshift=-2ex]\tikztotarget.west)
				-- (\tikztotarget)}] \\
			& \Hm^2(G, A) . & {} & {}
		\end{tikzcd}\end{equation*}
	\end{enumerate}

	Morfisme penghubung $\delta^0$ dan $\delta^1$ didefinisikan seperti di
	atas. Keduanya kanonik, yakni kompatibel dengan morfisme antarbarisan
	eksak pendek.
\end{teorema}
% segment-id: o014.li2.chapter6.nonabelian-cohomology.pf001
\begin{bukti}
	Konstruksi terperinci $\delta^0$ dan $\delta^1$ di atas telah cukup
	menunjukkan bahwa keduanya kanonik; dalam keadaan (iii), telah ditunjukkan
	pula bahwa $\delta^0$ merupakan homomorfisme grup. Berikut ini kita
	periksa keeksakan. Selanjutnya, pandang $A$ sebagai subgrup $B$ dan
	hilangkan $u$.
	
	Dalam keadaan (i), keeksakan pada $\Hm^0(G,A)$ dan $\Hm^0(G,B)$ jelas.
	Pada $\Hm^0(G,C)$, misalkan $c\in C^G$. Menurut definisi, $c$ dipetakan ke
	titik pangkal $\Hm^1(G,A)$ jika dan hanya jika terdapat $\alpha\in A$ dan
	$b\in B$ sedemikian sehingga
% segment-id: o014.li2.chapter6.nonabelian-cohomology.q019
	\[ v(b) = c, \quad \forall g \in G,\; b^{-1} \cdot {}^g b = \alpha^{-1} \cdot {}^g \alpha. \]
	Namun, ini ekuivalen dengan adanya $\alpha\in A$ dan $b\in B$ sedemikian
	sehingga $v(b)=c$ dan $b\alpha^{-1}\in B^G$, yang selanjutnya ekuivalen
	dengan $c\in\Image(\Hm^0(u))$.
	
	Pada $\Hm^1(G,A)$, misalkan $a\in Z^1(G,A)$. Kelas
	$[a]\in\Hm^1(G,A)$ dipetakan ke titik pangkal jika dan hanya jika terdapat
	$b\in B$ sedemikian sehingga
	$b^{-1}\cdot a(g)\cdot{}^gb=1$ selalu berlaku. Jika persamaan ini berlaku,
	ambillah $c:=v(b^{-1})$. Menerapkan $v$ pada
	${}^gb^{-1}=b^{-1}\cdot a(g)$ menunjukkan bahwa $c\in C^G$ dan
	$[a]=\delta^0([c])$. Argumen sebaliknya serupa.
	
	Dalam keadaan (ii), yang perlu dibuktikan ialah keeksakan pada
	$\Hm^1(G,B)$. Misalkan $b\in Z^1(G,B)$. Kelas $[b]$ dipetakan ke titik
	pangkal jika dan hanya jika terdapat $\beta\in B$ sedemikian sehingga
	$\beta^{-1}\cdot b(g)\cdot{}^g\beta\in A$ selalu berlaku. Jelas bahwa ini
	sama dengan mengatakan bahwa, hingga relasi $\sim$, kelas tersebut berasal
	dari $Z^1(G,A)$.
	
	Tinjau keadaan (iii). Misalkan $c\in Z^1(G,C)$. Kelas $[c]$ dipetakan ke
	titik pangkal jika dan hanya jika terdapat pemetaan $b:G\to B$ dan
	$\alpha:G\to A$ sedemikian sehingga $v\circ b=c$ dan
% segment-id: o014.li2.chapter6.nonabelian-cohomology.q020
	\[ b(g_1) \cdot {}^{g_1} b(g_2) = \alpha(g_1) \cdot {}^{g_1} \alpha(g_2) \cdot \alpha(g_1 g_2)^{-1} \cdot  b(g_1 g_2). \]
	Karena $A\subset Z_B$, jika $\alpha^{-1}b:G\to B$ didefinisikan dengan
	perkalian titik demi titik, persamaan di atas juga ekuivalen dengan
% segment-id: o014.li2.chapter6.nonabelian-cohomology.q021
	\[ (\alpha^{-1} b)(g_1) \cdot {}^{g_1} (\alpha^{-1} b)(g_2) = (\alpha^{-1} b)(g_1 g_2), \]
	yakni $\alpha^{-1}b\in Z^1(G,B)$. Jelas bahwa ini ekuivalen dengan
	pernyataan bahwa $[c]$ berasal dari $\Hm^1(G,B)$. Terbukti.
\end{bukti}

% segment-id: o014.li2.chapter6.nonabelian-cohomology.p019
Jika $A$ dan $B$ juga disyaratkan abelian, maka
$C\simeq B/u(A)$ merupakan grup abelian, sehingga kita kembali kepada versi
yang lazim dengan koefisien grup abelian.

% segment-id: o014.li2.chapter6.nonabelian-cohomology.p020
Perhatikan bahwa barisan eksak dalam
Teorema~\ref{prop:nonabelian-long} (i) mendeskripsikan prabayang titik pangkal
oleh $\Hm^1(u)$, tetapi secara umum tidak dapat menentukan apakah dua unsur
$\Hm^1(G,A)$ memiliki citra yang sama oleh $\Hm^1(u)$. Konstruksi
``puntiran'' yang diperkenalkan dalam latihan bab ini membantu menangani
masalah tersebut.

% segment-id: o014.li2.chapter6.nonabelian-cohomology.e001
\begin{contoh}
	Misalkan $B$ grup-$G$ dan $A$ subgrup yang stabil terhadap aksi $G$.
	Himpunan $B/A$ dilengkapi titik pangkal $1\cdot A$ dan aksi kiri $G$.
	Dengan demikian diperoleh barisan eksak pendek dalam
	$G\dcate{Set}_\bullet$
% segment-id: o014.li2.chapter6.nonabelian-cohomology.q022
	\[ 1 \to A \xrightarrow{\text{inklusi}} B \xrightarrow{\text{kuosien}} B/A \to 1. \]
	Syarat \eqref{eqn:pointed-ses-condition} jelas berlaku. Maka
	Teorema~\ref{prop:nonabelian-long} (i) memberikan barisan eksak dalam
	$\cate{Set}_\bullet$
% segment-id: o014.li2.chapter6.nonabelian-cohomology.q023
	\[ 1 \to A^G \to B^G \to (B/A)^G \xrightarrow{\delta^0} \Hm^1(G, A) \to \Hm^1(G, B). \]
	
	Pemetaan yang jelas $B^G/A^G\hookrightarrow(B/A)^G$ injektif, tetapi
	sering kali tidak bijektif: koset yang stabil di bawah aksi $G$ belum tentu
	memuat titik tetap-$G$; latihan akan memberikan contohnya. Meskipun
	demikian, barisan eksak tersebut mengidentifikasi obstruksi kohomologisnya.
	Untuk morfisme $f$ dalam $\cate{Set}_\bullet$, definisikan $\Ker(f)$
	sebagai prabayang titik pangkal. Maka
% segment-id: o014.li2.chapter6.nonabelian-cohomology.q024
	\begin{align*}
		\text{koset}\; bA \in (B/A)^G \;\text{berasal dari}\; B^G & \iff bA \in \Ker(\delta_0), \\
		(B/A)^G = B^G/A^G & \iff \Ker[\Hm^1(G, A) \to \Hm^1(G, B)] = 1.
	\end{align*}
\end{contoh}

% segment-id: o014.li2.chapter6.nonabelian-cohomology.p021
Bagian ini ditutup dengan versi nonabelian Lema Shapiro
(Teorema~\ref{prop:Shapiro}). Sebagai persiapan, terlebih dahulu definisikan
versi nonabelian modul terinduksi melalui ruang pemetaan; bandingkan
Lema~\ref{prop:Ind-as-mappings}. Konstruksi ini berlaku untuk objek
$H\dcate{Grp}$ atau $H\dcate{Set}_\bullet$, dengan $H$ subgrup dari $G$;
aksi grup dituliskan $(h,a)\mapsto{}^ha$.

% segment-id: o014.li2.chapter6.nonabelian-cohomology.d005
\begin{definisi}
	Misalkan $H$ subgrup $G$ dan $A$ objek $H\dcate{Grp}$ (atau
	$H\dcate{Set}_\bullet$). Definisikan objek $\Ind^G_H(A)$ dari
	$G\dcate{Grp}$ (atau $G\dcate{Set}_\bullet$) sebagai berikut. Sebagai
	objek $G\dcate{Set}$, objek tersebut ialah
% segment-id: o014.li2.chapter6.nonabelian-cohomology.q025
	\begin{gather*}
		\Ind^G_H(A) := \left\{\begin{array}{r|l}
			f: G \to A & \forall h \in H, \; \forall x \in G \\
			& f(hx) = {}^h f(x).
		\end{array}\right\}, \\
		({}^g f)(x) := f(xg), \quad f \in \Ind^G_H(A), \; x, g \in G,
	\end{gather*}
	dengan operasi grup (atau titik pangkal) didefinisikan titik demi titik pada
	pemetaan (atau oleh pemetaan konstan ke titik pangkal). Kita menyebut
	$A\mapsto\Ind^G_H(A)$ sebagai induksi.
\end{definisi}

% segment-id: o014.li2.chapter6.nonabelian-cohomology.p022
Konstruksi induksi tersebut menghasilkan funktor-funktor yang kompatibel,
$H\dcate{Grp}\to G\dcate{Grp}$ dan
$H\dcate{Set}_\bullet\to G\dcate{Set}_\bullet$.

% segment-id: o014.li2.chapter6.nonabelian-cohomology.th002
\begin{teorema}\label{prop:Shapiro-noncommutative}
	\index{Lema Shapiro}
	Misalkan $H$ subgrup $G$ dan $A$ objek $H\dcate{Set}_\bullet$ atau
	$H\dcate{Grp}$. Terdapat isomorfisme kanonik dalam $\cate{Set}_\bullet$
% segment-id: o014.li2.chapter6.nonabelian-cohomology.q026
	\[ \Hm^i\left(G, \Ind^G_H(A)\right) \rightiso \Hm^i(H, A), \]
	dengan $i=0$ jika $A$ berasal dari $H\dcate{Set}_\bullet$, atau $i=0,1$
	jika $A$ berasal dari $H\dcate{Grp}$. Pemetaannya adalah sebagai berikut.
% segment-id: o014.li2.chapter6.nonabelian-cohomology.l005
	\begin{description}
% segment-id: o014.li2.chapter6.nonabelian-cohomology.i010
		\item[($i=0$)] Ambil pemetaan $f\mapsto f(1_G)$ yang menentukan
		$Z^0\left(G,\Ind^G_H(A)\right)=\Ind^G_H(A)^G\to A^H=Z^0(H,A)$.
% segment-id: o014.li2.chapter6.nonabelian-cohomology.i011
		\item[($i=1$)] Ambil pemetaan
		$Z^1\left(G,\Ind^G_H(A)\right)\to Z^1(H,A)$ yang memetakan
		$c:G\to\Ind^G_H(A)$ ke
% segment-id: o014.li2.chapter6.nonabelian-cohomology.q027
		\begin{align*}
			(c|_H)(\cdot)(1_G): H & \to A \\
			h & \mapsto c(h)(1_G).
		\end{align*}
	\end{description}
	Keduanya menginduksi morfisme pada tingkat $\Hm^i$. Jika $A$ berasal dari
	$H\dcate{Grp}$, isomorfisme untuk $i=0$ juga merupakan isomorfisme grup.
\end{teorema}
% segment-id: o014.li2.chapter6.nonabelian-cohomology.pf002
\begin{bukti}
	Argumen rutin menunjukkan bahwa pemetaan dalam kedua kasus terdefinisi
	dengan baik dan mempertahankan titik pangkal atau struktur grup; perincian
	diserahkan kepada pembaca.
	
	Untuk $i=0$, mudah dilihat bahwa
	$\Ind^G_H(A)^G\to A^H$ memang bijektif dan memberikan isomorfisme grup
	jika $A$ grup-$H$. Berikut ini kita berfokus pada kasus ketika $A$ objek
	$H\dcate{Grp}$ dan $i=1$.
	
	Misalkan $r,s\in Z^1(G,\Ind^G_H(A))$ memiliki citra yang sama dalam
	$\Hm^1(H,A)$. Kita akan membuktikan $r\sim s$. Asumsi ini sama artinya
	dengan adanya $\alpha\in A$ sedemikian sehingga
% segment-id: o014.li2.chapter6.nonabelian-cohomology.q028
	\[ r(h)(1_G) = \alpha^{-1} \cdot s(h)(1_G) \cdot {}^h \alpha \]
	untuk setiap $h\in H$. Terdapat $\zeta\in\Ind^G_H(A)$ dengan
	$\zeta(1_G)=\alpha$. Dengan menyesuaikan $s$ oleh $\zeta$ dalam kelas
	ekuivalensinya, dapat dipastikan bahwa $r(h)(1_G)=s(h)(1_G)$ selalu berlaku.
	
	Definisi $Z^1(G,\Ind^G_H(A))$ menyiratkan kesamaan dalam $A$
% segment-id: o014.li2.chapter6.nonabelian-cohomology.q029
	\begin{gather*}
		r(g_1 g_2)(x) =r(g_1)(x) \cdot r(g_2)(xg_1), \quad s(g_1 g_2)(x) = s(g_1)(x) \cdot s(g_2)(xg_1)
	\end{gather*}
	untuk semua $x,g_1,g_2\in G$. Di satu pihak, dengan memasukkan
	$g_1=g\in G$ dan $g_2=g^{-1}x^{-1}$ lalu menyusun ulang, diperoleh
% segment-id: o014.li2.chapter6.nonabelian-cohomology.q030
	\begin{equation}\label{eqn:Shapiro-noncommutative-aux0}
		\begin{aligned}
			r(g)(x) & = r(x^{-1})(x) \cdot r(g^{-1} x^{-1})(xg)^{-1}, \\
			s(g)(x) & = s(x^{-1})(x) \cdot s(g^{-1} x^{-1})(xg)^{-1}.
		\end{aligned}
	\end{equation}
	Di pihak lain, dengan memasukkan $g_1=x^{-1}$ dan $g_2=h\in H$, diperoleh
% segment-id: o014.li2.chapter6.nonabelian-cohomology.q031
	\begin{equation}\label{eqn:Shapiro-noncommutative-aux1}
		s(x^{-1} h)(x) \cdot r(x^{-1} h)(x)^{-1} = s(x^{-1})(x) \cdot r(x^{-1})(x)^{-1}.
	\end{equation}
	
	Untuk setiap $x\in G$, definisikan
	$t(x):=s(x^{-1})(x)\cdot r(x^{-1})(x)^{-1}$. Kita menyatakan bahwa
	$t:G\to A$ termasuk dalam $\Ind^G_H(A)$. Memang, jika $h\in H$ dan
	$x\in G$, maka
% segment-id: o014.li2.chapter6.nonabelian-cohomology.q032
	\begin{multline*}
		t(hx) = s(x^{-1} h^{-1})(hx) \cdot r(x^{-1} h^{-1})(hx)^{-1} \\
		= {}^h \left( s(x^{-1}h^{-1})(x) \cdot r(x^{-1} h^{-1})(x)^{-1} \right) \\
		\xlongequal{\text{\eqref{eqn:Shapiro-noncommutative-aux1}}} {}^h \left( s(x^{-1})(x) \cdot r(x^{-1})(x)^{-1} \right) = {}^h t(x).
	\end{multline*}

	Untuk membuktikan $r\sim s$, cukup diperiksa identitas dalam
	$\Ind^G_H(A)$, yaitu $t^{-1}\cdot s(g)\cdot{}^gt=r(g)$. Nilai ruas kiri
	pada $x\in G$ adalah
% segment-id: o014.li2.chapter6.nonabelian-cohomology.q033
	\[ r(x^{-1})(x) \cdot s(x^{-1})(x)^{-1} \cdot s(g)(x) \cdot s(g^{-1} x^{-1})(xg) \cdot r(g^{-1} x^{-1})(xg)^{-1}; \]
	dengan menguraikan $s(g)(x)$ dan $r(g)(x)$ menggunakan
	\eqref{eqn:Shapiro-noncommutative-aux0}, jelas bahwa ekspresi ini sama
	dengan $r(g)(x)$. Injektivitas terbukti.
	
	Terakhir, buktikan surjektivitas. Misalkan $s^\flat\in Z^1(H,A)$.
	Tuliskan $\pi$ untuk pemetaan kuosien $G\twoheadrightarrow H\backslash G$,
	dan pilih sembarang pemetaan $\sigma:H\backslash G\to G$ dengan
	$\pi\sigma=\identity$. Tanpa mengurangi keumuman, dapat diandaikan
	$\sigma(H\cdot1_G)=1_G$. Setiap $g\in G$ memiliki representasi tunggal
% segment-id: o014.li2.chapter6.nonabelian-cohomology.q034
	\[ g = \tau(g) \sigma(Hg), \quad \tau(g) \in H. \]
	
	Perhatikan bahwa $\tau(1_G)=1_G$. Untuk setiap $g,x\in G$, definisikan
% segment-id: o014.li2.chapter6.nonabelian-cohomology.q035
	\[ s(g)(x) := {}^{\tau(x)} s^\flat\left( \tau(\sigma(Hx)g) \right) \; \in A. \]
	Jika $h\in H$, dari $\tau(hx)=h\tau(x)$ segera diperoleh
	$s(g)(hx)={}^hs(g)(x)$, sehingga $s(g)\in\Ind^G_H(A)$. Langkah berikutnya
	ialah menunjukkan $s\in Z^1(G,\Ind^G_H(A))$. Pembaca terlebih dahulu
	diminta memeriksa
% segment-id: o014.li2.chapter6.nonabelian-cohomology.q036
	\begin{equation}\label{eqn:Shapiro-noncommutative-aux2}
		\tau(uv) = \tau(u) \tau(\sigma(Hu)v), \quad u, v \in G.
	\end{equation}
	
	Tujuannya ialah membuktikan
	$s(g_1g_2)(x)=s(g_1)(x)\cdot s(g_2)(xg_1)$ untuk semua
	$g_1,g_2,x\in G$. Berdasarkan definisi $Z^1(H,A)$, ruas kiri sama dengan
% segment-id: o014.li2.chapter6.nonabelian-cohomology.q037
	\begin{multline*}
		{}^{\tau(x)} s^\flat\left( \tau(\sigma(Hx)g_1 g_2) \right) \xlongequal{\text{\eqref{eqn:Shapiro-noncommutative-aux2}}} {}^{\tau(x)} s^\flat\left( \tau(\sigma(Hx)g_1) \tau(\sigma(Hx g_1)g_2) \right) \\
		= {}^{\tau(x)} s^\flat\left( \tau(\sigma(Hx)g_1) \right) \cdot {}^{\tau(x)\tau(\sigma(Hx)g_1)} s^\flat\left( \tau(\sigma(Hxg_1)g_2) \right) \\
		\xlongequal{\text{\eqref{eqn:Shapiro-noncommutative-aux2}}} {}^{\tau(x)} s^\flat\left( \tau(\sigma(Hx)g_1) \right) \cdot {}^{\tau(xg_1)} s^\flat\left( \tau(\sigma(Hxg_1)g_2) \right),
	\end{multline*}
yang merupakan ruas kanan.

	Selain itu, jika $h\in H$, maka $s(h)(1_G)=s^\flat(h)$, sehingga $s$
	dipetakan ke $s^\flat$. Surjektivitas terbukti.
\end{bukti}

% segment-id: o014.li2.chapter6.nonabelian-cohomology.p023
Jelas bahwa isomorfisme dalam
Teorema~\ref{prop:Shapiro-noncommutative}, jika $A$ grup-$G$ abelian,
mereduksi kepada kasus khusus Proposisi~\ref{prop:Shapiro-std-coh}.
Selain itu, perhitungan langsung menunjukkan bahwa isomorfisme ini
kompatibel dengan morfisme penghubung dalam
Teorema~\ref{prop:nonabelian-long}; perinciannya diserahkan sebagai latihan.

% segment-id: o014.li2.chapter6.nonabelian-cohomology.r002
\begin{catatan}[Kohomologi nonabelian grup profinit]\label{rem:pro-finite-coh-nonabelian}
	\index{grup-G@grup-$G$!mulus (smooth)}
	Tinjau grup profinit $G$. Objek $A$ dari $G\dcate{Set}_\bullet$ atau
	$G\dcate{Grp}$ disebut \emph{mulus} jika
	$A=\bigcup_KA^K$, dengan $K$ berkisar atas subgrup normal terbuka dari $G$;
	dalam keadaan ini, $A$ diberi topologi diskret. Karena argumen dalam
	bagian ini semuanya didasarkan pada analogi dengan kompleks kokrantai
	standar, semua hasil sebelumnya berlaku bagi grup profinit $G$ dan $A$
	yang mulus. Perbedaannya hanya bahwa $c\in Z^i(G,A)$, sebagai pemetaan dari
	$G$ ke $A$, harus kontinu, yakni lokal konstan. Dengan demikian,
	$\Hm^i(G,A)$ dan berbagai morfisme kanonik dapat didefinisikan untuk
	$i=0,1$, dan
% segment-id: o014.li2.chapter6.nonabelian-cohomology.q038
	\begin{align*}
		Z^i(G, A) & = \varinjlim_K Z^i(G/K, A^K), \\
		\Hm^i(G, A) & = \varinjlim_K \Hm^i(G/K, A^K), \quad i = 0, 1.
	\end{align*}

	Selain itu, definisi $\Ind^G_H(A)$ harus mensyaratkan bahwa $f:G\to A$
	lokal konstan. Bukti Teorema~\ref{prop:Shapiro-noncommutative} juga
	memerlukan Lema~\ref{prop:profinite-section} untuk memastikan bahwa
	$\zeta:G\to A$ dalam bagian injektivitas dapat dipilih lokal konstan dan
	$\sigma:H\backslash G\to G$ dalam bagian surjektivitas kontinu.
	
	Terakhir, Catatan~\ref{rem:Z1-vs-section} tetap berlaku bagi grup profinit
	$G$: $Z^1(G,A)$ berkorespondensi dengan himpunan penampang kontinu dari
	homomorfisme proyeksi $\pi:A\rtimes G\to G$, sedangkan $\Hm^1(G,A)$ ialah
	hasil bagi himpunan penampang kontinu oleh aksi konjugasi $A$.
\end{catatan}
